\section{Majorization between density matrices}

\begin{definition}[Majorization of Hermitian matrices]
    \label{def:density_matrix_majorization}
    \lean{Matrix.IsHermitian.Majorizes}
    \leanok
    \uses{def:vector_majorization}
    Let $\rho,\sigma\in M_d(\C)$ be Hermitian. Write $\rho\succ\sigma$
    when the eigenvalue vector of $\sigma$ is majorized by the eigenvalue
    vector of $\rho$. Eigenvalues are counted with their multiplicities;
    majorization compares their decreasing rearrangements.
\end{definition}

\begin{definition}[Mixture of unitary conjugations]
    \label{def:density_unitary_mixture}
    \lean{TNLean.IsRandomUnitaryImage}
    \leanok
    A matrix $\sigma$ is a finite probability mixture of unitary
    conjugations of $\rho$ if there are finitely many unitary matrices
    $U_i\in M_d(\C)$ and real numbers $p_i\geq0$ with $\sum_i p_i=1$ such that
    \begin{align}
        \sigma&=\sum_i p_i U_i\rho U_i^\dagger.
        \label{eq:density_unitary_mixture}
    \end{align}
    Zero coefficients and repeated unitaries are allowed.
\end{definition}

\begin{lemma}[Squared entries of a unitary matrix]
    \label{lem:density_unitary_stochastic}
    \lean{TNLean.unitary_normSq_mem_doublyStochastic}
    \leanok
    If $U\in M_d(\C)$ is unitary, the matrix with entries $|U_{kl}|^2$
    is doubly stochastic.
\end{lemma}
\begin{proof}\leanok
    Its entries are nonnegative. The diagonal entries of $UU^\dagger=I$
    and $U^\dagger U=I$ give its row and column sums, respectively.
\end{proof}

\begin{theorem}[Majorization and mixtures of unitary conjugations]
    \label{thm:density_matrix_majorization}
    \lean{Matrix.IsHermitian.majorizes_iff_trace_convex,
        Matrix.IsHermitian.Majorizes.isRandomUnitaryImage,
        Matrix.IsHermitian.majorizes_of_randomUnitaryImage,
        Matrix.IsHermitian.majorizes_iff_randomUnitaryImage}
    \leanok
    Let $\rho,\sigma\in M_d(\C)$ be Hermitian. The following are equivalent:
    \begin{enumerate}
        \item $\rho\succ\sigma$;
        \item $\operatorname{Re}\tr f(\sigma)\leq\operatorname{Re}\tr f(\rho)$
            for every convex function $f:\R\to\R$;
        \item $\sigma$ is a finite probability mixture of unitary
            conjugations of $\rho$.
    \end{enumerate}
    Here $f(\rho)$ and $f(\sigma)$ are defined by spectral functional calculus.
    In particular the equivalence holds for every pair of density matrices,
    as in \cite[Theorem~8.7 and Equation~(8.21)]{Wolf2012Quantum}.
\end{theorem}
\begin{proof}\leanok
    \uses{def:density_matrix_majorization, def:density_unitary_mixture,
        lem:density_unitary_stochastic, thm:vector_majorization_equivalences,
        thm:birkhoff_doubly_stochastic}
    Choose spectral decompositions
    $\rho=V_\rho\operatorname{diag}(\lambda^\rho)V_\rho^\dagger$ and
    $\sigma=V_\sigma\operatorname{diag}(\lambda^\sigma)V_\sigma^\dagger$.
    Every function on the finite spectra admits functional calculus, and
    $\operatorname{Re}\tr f(\rho)=\sum_j f(\lambda^\rho_j)$, with the
    analogous identity for $\sigma$. The convex trace characterization
    therefore follows from vector majorization.

    If $\rho\succ\sigma$, vector majorization gives a doubly stochastic
    matrix $M$ with $\lambda^\sigma=M\lambda^\rho$.
    Write $M=\sum_\tau p_\tau P_\tau$ by Birkhoff's theorem, with
    $p_\tau\geq0$ and $\sum_\tau p_\tau=1$.
    Conjugating the resulting diagonal-matrix identity gives
    (\ref{eq:density_unitary_mixture}) with
    $U_\tau=V_\sigma P_\tau V_\rho^\dagger$.
    These are unitary matrices, so this constructs the required mixture.

    Conversely, suppose (\ref{eq:density_unitary_mixture}) holds and set
    $W_i=V_\sigma^\dagger U_iV_\rho$.
    Conjugating the mixture by $V_\sigma^\dagger$ and comparing diagonal
    entries gives
    \begin{align}
        \lambda^\sigma_k&=\sum_l M_{kl}\lambda^\rho_l,
        &M_{kl}&=\sum_i p_i |(W_i)_{kl}|^2.
        \label{eq:density_mixture_stochastic}
    \end{align}
    Each squared-entry matrix is doubly stochastic, and their probability
    mixture $M$ is doubly stochastic. Vector majorization completes the proof.
\end{proof}

\begin{lemma}[Positive dimension of a density matrix]
    \label{lem:density_dimension_positive}
    \lean{Matrix.dimension_pos_of_trace_one}
    \leanok
    A matrix of trace $1$ has positive dimension. In particular there are
    no density matrices in dimension zero.
\end{lemma}
\begin{proof}\leanok
    The trace of the empty matrix is $0$.
\end{proof}
\endinput
